%%%%%%%%%%%%%%%%%%%%%%%%%%%%%%%%%%%%%%%%%%%%%%%%%%%%%%%%%%%%%%%%%%%%%%%%
% Uni Duesseldorf
% Lehrstuhl fuer Datenbanken und Informationssysteme
% Vorlage fuer Bachelor-/Masterarbeiten
% Optimiert fuer den Original-Latex-Kompiler LATEX.EXE (LaTeX=>PS=>PDF)
%%%%%%%%%%%%%%%%%%%%%%%%%%%%%%%%%%%%%%%%%%%%%%%%%%%%%%%%%%%%%%%%%%%%%%%%
% Ueberarbeitung für pdflatex (LaTeX=>PDF)
%%%%%%%%%%%%%%%%%%%%%%%%%%%%%%%%%%%%%%%%%%%%%%%%%%%%%%%%%%%%%%%%%%%%%%%%
% Vorlage Changelog:
% 10.09.2015 (Matthias Liebeck): Nummerierung des Inhaltsverzeichnis nun römisch, Beispiel für einen Anhang eingebaut, \raggedbottom hinter sections eingefügt
% 11.07.2018 (Matthias Liebeck): Ersetzung des Bibliographiestils, Einsatz von Biber
% 04.09.2018 (Matthias Liebeck):
%   * Bibtex: unnötige Bibtexfelder beim Rendern ausblenden (thx @ Markus Brenneis)
%   * ngerman: "et al." im BibTeX für drei oder mehr Autoren
%   * Neuer Befehl \sectionforcestartright: Sections immer rechts beginnen (thx @ Philipp Grawe)
%   * ngerman: Deutsche Anführungszeichen im Literaturverzeichnis (thx @ Markus Brenneis)
%   * ngerman: Deutsche Anführungszeichen im Literaturverzeichnis (thx @ Markus Brenneis)
% 16.10.2018 (Matthias Liebeck): Zwei fixes an \sectionforcestartright (thx @ Markus Brenneis)
%%%%%%%%%%%%%%%%%%%%%%%%%%%%%%%%%%%%%%%%%%%%%%%%%%%%%%%%%%%%%%%%%%%%%%%%
%%%% BEGINN EINSTELLUNG FUER DIE ARBEIT. UNBEDINGT ERFORDERLICH! %%%%%%%
%%%%%%%%%%%%%%%%%%%%%%%%%%%%%%%%%%%%%%%%%%%%%%%%%%%%%%%%%%%%%%%%%%%%%%%%
% Geben Sie Ihren Namen hier an:
\newcommand{\bearbeiter}{Felix Michels}

% Geben Sie hier den Titel Ihrer Arbeit an:
\newcommand{\titel}{Adversarial Attacks on Capsule Networks}

% Geben Sie das Datum des Beginns und Ende der Bachelorarbeit ein:
\newcommand{\beginndatum}{08. April 2019}
\newcommand{\abgabedatum}{???}

% Geben Sie die Namen des Erst- und Zweitgutachters an:
\newcommand{\erstgutachter}{Prof. Dr.~Stefan Harmeling}
\newcommand{\zweitgutachter}{Prof. Dr.~Stefan Conrad}

% Falls Sie die Arbeit zweiseitig ausdrucken wollen,
% benutzen Sie die folgende Zeile mit
% \AN fuer zweiseitigen Druck
% \AUS fuer einseitigen Druck
\newcommand{\zweiseitig}{\AN}

% Falls Sections immer rechts beginnen sollen. Gerade für Masterarbeiten
% interessant. Bei kurzen Bachelorarbeiten eher weniger zu verwenden.
\newcommand{\sectionforcestartright}{\AUS}
%\newcommand{\sectionforcestartright}{\AN}

% Falls die Arbeit in englischer Sprache verfasst 
% werden soll, dann benutzen Sie die folgende Zeile mit
% englisch fuer englische Sprache
% deutsch fuer deutsche Sprache
\newcommand{\sprache}{englisch}

% Hier wird eingestellt, ob es sich bei der Arbeit um eine Bachelor- 
% oder Masterarbeit handelt (unpassendes auskommentieren!):
\newcommand{\arbeit}{Bachelorarbeit}
%~ \newcommand{\arbeit}{Masterarbeit}


%%%%%%%%%%%%%%%%%%%%%%%%%%%%%%%%%%%%%%%%%%%%%%%%%%%%%%%%%%%%%%%%%%%%%%%%
%%%% ENDE EINSTELLUNGEN %%%%%%%%%%%%%%%%%%%%%%%%%%%%%%%%%%%%%%%%%%%%%%%%
%%%%%%%%%%%%%%%%%%%%%%%%%%%%%%%%%%%%%%%%%%%%%%%%%%%%%%%%%%%%%%%%%%%%%%%%

% Die folgende Zeile NICHT EDITIEREN oder loeschen


%%%%%%%%%%%%%%%%%%%%%%%%%%%%%%%%%%%%%%%%%%%%%%%%%%%%%%%%%%%
% Obere Titelmakros. Editieren Sie diese Datei nur, wenn
% Sie sich ABSOLUT sicher sind, was Sie da tun!!!
% (Z.B. zum Abaendern der BA-Vorlage in eine MA-Vorlage)
% Uni Duesseldorf
% Lehrstuhl fuer Datenbanken und Informationssysteme
% Version 2.2 - 2.3.2010
%%%%%%%%%%%%%%%%%%%%%%%%%%%%%%%%%%%%%%%%%%%%%%%%%%%%%%%%%%%
\newcommand{\AN}{twoside}
\newcommand{\AUS}{}
%\newcommand{\englisch}{}
%\newcommand{\deutsch}{\usepackage[german]{babel}}

%% Die folgenden auskommentierten Optionen dienen der automatischen
%% Erkennung des Latex-Kompilers und dem Setzen der davon abhängigen
%% Einstellungen. Bei Problem z.B. mit dem Einbinden von verschiedenen
%% Grafiktypen bei Verwendung von PdfLatex oder Latex, einfach die
%% verschiedenen \usepackage(s) ausprobieren. (Mit diesen Einstellungen
%% funktionierte diese Vorlage bei der Verwenundg von latex.exe als
%% Kompiler bei den meisten Studierenden.)

%\newif\ifpdf \ifx\pdfoutput\undefined
%\pdffalse % we are not running pdflatex
%\else
%\pdfoutput=1 % we are running pdflatex
%\pdfcompresslevel=9 % compression level for text and image;
%\pdftrue \fi

\documentclass[11pt,a4paper, \zweiseitig]{article}



%\usepackage[iso]{umlaute}
\usepackage[utf8]{inputenc}
\usepackage{palatino} % palatino Schriftart
%\usepackage{makeidx} % um ein Index zu erstellen
\usepackage[nottoc]{tocbibind}
\usepackage[T1]{fontenc} %fuer richtige Trennung bei Umlauten
\usepackage{fancybox} % fuer die Rahmen
\usepackage{shortvrb}
\usepackage{url}
\usepackage{xcolor}
\usepackage[colorlinks,citecolor=blue,linkcolor=black]{hyperref} %anklickbares Inhaltsverzeichnis
\usepackage[style=authoryear,natbib=true,backend=biber,mincitenames=1,maxcitenames=2,maxbibnames=99,uniquelist=false,dashed=false]{biblatex}

%% Extra packages
\usepackage[colorinlistoftodos,obeyFinal]{todonotes}
\usepackage{graphicx}
\usepackage{booktabs}
\usepackage{multirow}

%% Math
\usepackage{amsmath}
\usepackage{amssymb}

\DeclareMathOperator*{\minimize}{minimize}

% https://tex.stackexchange.com/a/334703/8850
\AtEveryBibitem{%
  \clearfield{issn}
  \clearfield{isbn}
  \clearfield{doi}
  \clearfield{location}
  \clearlist{location}
  \clearlist{address}

  \ifentrytype{online}{}{% Remove url except for @online
    \clearfield{url}
  }
}

% Falls es bei \citet ein Komma zwischen Name und Jahr gibt:
% https://tex.stackexchange.com/questions/312539/unwanted-comma-between-author-and-year-using-citet-command
% (thx @ Markus Brenneis)
%\DeclareDelimFormat[cbx@textcite]{nameyeardelim}{\addspace}



\usepackage{ifthen}
\ifthenelse{\equal{\sprache}{deutsch}}{
 \usepackage[ngerman]{babel}
 % Bibtex u.a -> et al.
 \DefineBibliographyStrings{ngerman}{
   andothers = {{et\,al\adddot}},
 }
 \usepackage[autostyle, german=quotes]{csquotes} % Deutsche Anführungszeichen im Literaturverzeichnis (thx @ Markus Brenneis)
 }{}

\usepackage{a4wide} % ganze A4 Weite verwenden



%\ifpdf
%\usepackage[pdftex,xdvi]{graphicx}
%\usepackage{thumbpdf} %thumbs fuer Pdf
%\usepackage[pdfstartview=FitV]{hyperref} %anklickbares Inhaltsverzeichnis
%\else
%\usepackage[dvips,xdvi]{graphicx}
\usepackage{graphicx}

%\fi

\newcommand{\redt}[1] {
  \textcolor{red}{#1}}

\newcommand{\oranget}[1] {
  \textcolor{orange}{#1}}

\newcommand{\purplet}[1] {
  \textcolor{purple}{#1}}

%%%%%%%%%%%%%%%%%%%%%%% Massangaben fuer die Arbeit %%%%%%%%%%%%%%%
\setlength{\textwidth}{15cm}

\setlength{\oddsidemargin}{35mm}
\setlength{\evensidemargin}{25mm}

\addtolength{\oddsidemargin}{-1in}
\addtolength{\evensidemargin}{-1in}

\addbibresource{references.bib}

%\makeindex

\begin{document}

%\setcounter{secnumdepth}{4} %Nummerieren bis in die 4. Ebene
%\setcounter{tocdepth}{4} %Inhaltsverzeichnis bis zur 4. Ebene

\pagestyle{headings}

\sloppy % LaTeX ist dann nicht so streng mit der Silbentrennung
%~ \MakeShortVerb{\§}

\parindent0mm
\parskip0.5em


{
\textwidth170mm 
\oddsidemargin30mm 
\evensidemargin30mm 
\addtolength{\oddsidemargin}{-1in}
\addtolength{\evensidemargin}{-1in}

\parskip0pt plus2pt

% Die Raender muessen eventuell fuer jeden Drucker individuell eingestellt
% werden. Dazu sind die Werte fuer die Abstaende `\oben' und `\links' zu
% aendern, die von mir auf jeweils 0mm eingestellt wurden.

%\newlength{\links} \setlength{\links}{10mm}  % hier abzuaendern
%\addtolength{\oddsidemargin}{\links}
%\addtolength{\evensidemargin}{\links}

\begin{titlepage}
\vspace*{-1.5cm}
  \raisebox{17mm}{
    \begin{minipage}[t]{70mm}
      \begin{center}
        %\selectlanguage{german}
        {\Large INSTITUT FÜR INFORMATIK\\}
        {\normalsize
          Machine Learning\\
        }
        \vspace{3mm}
        {\small Universitätsstr. 1 \hspace{5ex} D--40225 Düsseldorf\\}
     \end{center}
    \end{minipage}
  }
  \hfill
  \includegraphics[width=130pt]{bilder/HHU_Logo}
  \vspace{14em}

% Titel
  \begin{center}
      	\baselineskip=55pt
    	\textbf{\huge \titel}
  	 	\baselineskip=0 pt
   \end{center}

  %\vspace{7em}

\vfill

% Autor
  \begin{center}
    \textbf{\Large
      \bearbeiter
    }
  \end{center}

  \vspace{35mm}
 
% Prüfungsordnungs-Angaben
  \begin{center}
    %\selectlanguage{german}
    
%%%%%%%%%%%%%%%%%%%%%%%%%%%%%%%%%%%%%%%%%%%%%%%%%%%%%%%%%%%%%%%%%%%%%%%%%
% Ja, richtig, hier kann die BA-Vorlage zur MA-Vorlage gemacht werden...
% (nicht mehr nötig!)
%%%%%%%%%%%%%%%%%%%%%%%%%%%%%%%%%%%%%%%%%%%%%%%%%%%%%%%%%%%%%%%%%%%%%%%%%
    {\Large \arbeit}

    \vspace{2em}

    \begin{tabular}[t]{ll}
      Beginn der Arbeit:& \beginndatum \\
      Abgabe der Arbeit:& \abgabedatum \\
      Gutachter:         & \erstgutachter \\
                         & \zweitgutachter \\
    \end{tabular}
  \end{center}

\end{titlepage}

}

%%%%%%%%%%%%%%%%%%%%%%%%%%%%%%%%%%%%%%%%%%%%%%%%%%%%%%%%%%%%%%%%%%%%%
\clearpage
\begin{titlepage}
  ~                % eine leere Seite hinter dem Deckblatt
\end{titlepage}
%%%%%%%%%%%%%%%%%%%%%%%%%%%%%%%%%%%%%%%%%%%%%%%%%%%%%%%%%%%%%%%%%%%%%
\clearpage
\begin{titlepage}
\vspace*{\fill}

\section*{Erklärung}

%%%%%%%%%%%%%%%%%%%%%%%%%%%%%%%%%%%%%%%%%%%%%%%%%%%%%%%%%%%
% Und hier ebenfalls ggf. BA durch MA ersetzen...
% (Auch nicht mehr nötig!)
%%%%%%%%%%%%%%%%%%%%%%%%%%%%%%%%%%%%%%%%%%%%%%%%%%%%%%%%%%%

Hiermit versichere ich, dass ich diese \arbeit{}
selbstständig verfasst habe. Ich habe dazu keine anderen als die
angegebenen Quellen und Hilfsmittel verwendet.

\vspace{25 mm}

\begin{tabular}{lc}
Düsseldorf, den \abgabedatum \hspace*{2cm} & \underline{\hspace{6cm}}\\
& \bearbeiter
\end{tabular}

\vspace*{\fill}
\end{titlepage}

%%%%%%%%%%%%%%%%%%%%%%%%%%%%%%%%%%%%%%%%%%%%%%%%%%%%%%%%%%%%%%%%%%%%%
% Leerseite bei zweiseitigem Druck
%%%%%%%%%%%%%%%%%%%%%%%%%%%%%%%%%%%%%%%%%%%%%%%%%%%%%%%%%%%%%%%%%%%%%

\ifthenelse{\equal{\zweiseitig}{twoside}}{\clearpage\begin{titlepage}
~\end{titlepage}}{}

%%%%%%%%%%%%%%%%%%%%%%%%%%%%%%%%%%%%%%%%%%%%%%%%%%%%%%%%%%%%%%%%%%%%%
\clearpage
\begin{titlepage}

%%% Die folgende Zeile nicht ändern!
\section*{\ifthenelse{\equal{\sprache}{deutsch}}{Zusammenfassung}{Abstract}}
%%% Zusammenfassung:

This bachelor thesis extensively evaluates the vulnerability of capsule network to various adversarial attacks.
Capsule networks are a form of neural network in which neurons are grouped in capsules whose activity vectors describe different properties of the same entity.
Recent work attributes them more robustness to adversarial attacks than other types of neural networks.

To investigate those conjectures we train four pairs of capsule networks and convolutional neural networks discriminatively on different image datasets.
Using multiple strong and diverse adversarial attack we compute adversarial examples for each combination of dataset and architecture.

By comparing the euclidean norms of the generated perturbations we find that capsule networks are vulnerable to adversarial attacks in a similar level as convolutional networks.
We discover in particular that capsule networks show no robustness to sufficiently strong white-box attacks.
Additionally we demonstrate limited transferability of the adversarial examples between the two types of architectures.

We study furthermore the structure of the adversarial perturbations and can observe that they seem to lie in a relatively low dimensional linear subspace in both cases.
Despite those similarities, using dimensionality reduction we detect inherent differences between the perturbations constructed for the different network types.

Finally we aim to identify the source of adversarial examples in the networks we trained.
Since capsule networks and convolutional partially use the same type of convolutional layer we focus on the disturbance on the activation of individual layers caused by adversarial attacks.
We find that susceptibility to attacks is not just ascribable to the first layers in a network but a result of deeper layers as well.
Specifically we find that the dynamic routing algorithm used in capsule networks does not establish capsule activation invariants to counteract adversarial attacks as previously conjectured.


%%%%%%%%%%%%%%%%%%%%%%%%%%%%%%%%%%%%%%%%%%%%%%%%
% Untere Titelmakros. Editieren Sie diese Datei nur, wenn Sie sich
% ABSOLUT sicher sind, was Sie da tun!!!
%%%%%%%%%%%%%%%%%%%%%%%%%%%%%%%%%%%%%%%%%%%%%%%
\vspace*{\fill}
\end{titlepage}

%%%%%%%%%%%%%%%%%%%%%%%%%%%%%%%%%%%%%%%%%%%%%%%%%%%%%%%%%%%%%%%%%%%%%
% Leerseite bei zweiseitigem Druck
%%%%%%%%%%%%%%%%%%%%%%%%%%%%%%%%%%%%%%%%%%%%%%%%%%%%%%%%%%%%%%%%%%%%%
\ifthenelse{\equal{\zweiseitig}{twoside}}
  {\clearpage\begin{titlepage}~\end{titlepage}}{}
%%%%%%%%%%%%%%%%%%%%%%%%%%%%%%%%%%%%%%%%%%%%%%%%%%%%%%%%%%%%%%%%%%%%%
\clearpage \setcounter{page}{1}
\pagenumbering{roman}
\setcounter{tocdepth}{2}
\tableofcontents

%\enlargethispage{\baselineskip}
\clearpage
%%%%%%%%%%%%%%%%%%%%%%%%%%%%%%%%%%%%%%%%%%%%%%%%%%%%%%%%%%%%%%%%%%%%%
% Leere Seite, falls Inhaltsverzeichnis mit ungerader Seitenzahl und 
% doppelseitiger Druck
%%%%%%%%%%%%%%%%%%%%%%%%%%%%%%%%%%%%%%%%%%%%%%%%%%%%%%%%%%%%%%%%%%%%%
\ifthenelse{ \( \equal{\zweiseitig}{twoside} \and \not \isodd{\value{page}} \)}
	{\pagebreak \thispagestyle{empty} \cleardoublepage}{\clearpage}


% Kapitel soll bei doppelseitigem Druck immer auf der rechten (ungeraden) Seite anfangen (thx @ Philipp Grawe)
% https://tex.stackexchange.com/a/223387
\ifthenelse{\( \equal{\sectionforcestartright}{twoside} \)}
 {\let\oldsection\section % Store \section in \oldsection
   \renewcommand{\section}{\cleardoublepage\oldsection}}
 {}
\pagenumbering{arabic}
\setcounter{page}{1}

%%%%%%%%%%%%%%%%%%%%%%%%%%%%%%%%%%%%%%%%%%%%%%%%%%%%%%%%%%%%%%%%%%%%%%%%
%%%% BEGINN TEXTTEIL %%%%%%%%%%%%%%%%%%%%%%%%%%%%%%%%%%%%%%%%%%%%%%%%%%%
%%%%%%%%%%%%%%%%%%%%%%%%%%%%%%%%%%%%%%%%%%%%%%%%%%%%%%%%%%%%%%%%%%%%%%%%

%%%%%%%%%%%%%%%%%%%%%%%%%%%%%%%%%%%%%%%%%%%%%%%%%%%%%%%%%%%%%%%%%%%%%%%%
% Text entweder direkt hier hinein schreiben oder, im Sinne der
% besseren Uebersichtlich- und Bearbeitbarkeit mittels \input die
% einzelnen Textteile hier einbinden.
%%%%%%%%%%%%%%%%%%%%%%%%%%%%%%%%%%%%%%%%%%%%%%%%%%%%%%%%%%%%%%%%%%%%%%%%

\section{Introduction}
In this bachelor thesis I will show, that capsule networks are horrible
% !TeX root = ba.tex

\section{Capsule Networks}
\label{sec:capsules}

\todo{Rewrite this whole section}

The concept of vector capsules and the dynamic routing algorithm was proposed by \citet{capsules}. In an essence, neurons are grouped into vectors, so-called capsules. Each capsule vector is dedicated to a distinct abstract entity, i.e. a single object class in a classification setting. The norm of a capsule vector encodes the probability of the represented object being present in the input, while the vector orientation encodes the object's characteristics. Thus, CapsNets aim to develop dedicated representations that are distributed into multiple vectors in contrast to convolutional networks that utilize an entangled representation in a single vector at a given location. This allows the application of linear transformations directly to the representations of respective entities. Spatial relations, which can be implemented as a matrix product, can thus be modeled more efficiently.

CapsNets are organized in layers. Initially, the original CapsNet applies a convolutional layer. The resulting feature maps are then processed by the primary capsule layer. Internally, it applies a series of convolutional layers on its own, each yielding a spatial grid of capsules. Within all capsule layers the \emph{squashing} function serves as a vector-to-vector non-linearity that squashes each capsule vector length between $0$ and $1$ while leaving the orientation unaltered. Subsequently, convolutional or densely connected capsule layers can be applied. While the latter does not utilize weight sharing, convolutional capsule layers share the kernels over the spatial grid, as well as capsules from the previous layer. These layers effectively estimate output capsules based on respective input capsules. The dynamic routing algorithm determines each agreement between estimate and iteratively calculated output capsule. This is done by introducing a scalar factor, the so-called routing coefficient, for each connection between an estimate and respective output. Such an output is defined as the sum over all respective estimates, weighted by their routing coefficients. Theoretically, that means information flows where it is needed, both during forward and backpropagation. This non-parametric procedure supports the goal of capsules with clean dedicated representations. To improve results, an additional capsule may be used within the routing algorithm to serve as a dead end for information that may not be linked to known abstract capsule categories. This is also referred to as the \emph{none-of-the-above} category.

- History of capsule networks (Hinton, Sabour) \\
- Explain Dynamic Routing \\
- Mention other routing algorithms \\
- Parse tree \\
- Representation??? \\

Squashing:
\begin{equation}
\begin{aligned}
& v_j = \frac{\norm{s_j}}{1 + \norm{s_j}^2} s_j \\
& \sum_{i} c_{ij} \hat{u}_{j \vert i}, & & \hat{u}_{j \vert i} = W_{ij}u_j
\end{aligned}
\end{equation}

\algtext*{EndFor} % Remove end for text
\begin{algorithm}
\caption{Routing-by-agreement}
\label{alg:foobar}
\begin{algorithmic}[1]
	
\Procedure{Foo}{$x$}
	\State $\forall i,j: b_{ij} \gets 0$
	\For{$r$ iterations}
		\State $c_{ij} \gets \frac{\exp(b_{ij})}{\sum_{k}{\exp(b_{ik})}}$ \Comment{Softmax}
		\State $s_j \gets \sum_{i}{c_{ij}\hat{u}_{j \vert i}}$
		\State $v_j = \frac{\norm{s_j}}{1 + \norm{s_j}^2} s_j$ \Comment{Squash capsules}
		\State $b_{ij} \gets b_{ij} + \langle v_j, \hat{u}_{j \vert i} \rangle$
	\EndFor
\EndProcedure
	
\end{algorithmic}
\end{algorithm}
% !TeX root = ba.tex

\section{Adversarial Attacks}
\label{sec:attacks}

In general, adversarial attacks describe any techniques with the purpose to disturb the performance of machine learning models.
While adversarial attacks exist for a wide variety of applications like speech recognition or natural language processing, this work will concentrate on adversarial attacks on image classification tasks.
\emph{Adversarial examples} refer to modified images which are very close to some original image, but are classified differently by the model.
In our context they are in particular only mentioned in relationship to the original image.
The \emph{adversarial perturbation} is the difference between the adversarial example and the original.

There are various attack scenarios differing in the knowledge, goals and abilities of the attacker.
In \emph{white-box} attacks, all information about the model architecture, training images and learned parameters is available to the attacker, while in the \emph{black-box} scenario only the networks output can be observed.
In particular, gradients with respect to the input can be computed in the white-box setting.
These terms are in general not strictly defined and many intermediate stages exist, including limited access to the models output, information about the architecture but not the learned parameters or knowledge about the training images but not the architecture etc.

\emph{Targeted} and \emph{untargeted} describes the objective of the attacker.
With an untargeted attack an adversarial example is just classified as any other label than the true label, while an adversarial example from a targeted attack has to be classified as a given target label chosen by the attacker beforehand.
Targeted attacks are strictly more difficult for the attacker, but also more useful.
Furthermore, in some instances untargeted adversarial examples are easy to generate, or rather are near impossible to defend against, just because some classes may be conceptually very close to each other.
For example, differentiating between a Golden Retriever and a Labrador Retriever in ImageNet \citep{imagenet} can even be quite challenging for humans.

The last type of adversarial attacks that we will test are \emph{universal perturbations}.
Most attacks compute a different adversarial perturbation for each input image, but a universal perturbation aims to fool many images while using the same perturbation.
Non-universal perturbations are easier to find, but universal perturbations are simpler for the attacker to apply in a realistic attack scenario.
Additionally, the mere existence of universal perturbations reveals significant flaws in the used architecture, more so than that of normal adversarial examples.

An important problem to consider is the definition of \emph{closeness} to the original image. Possibilities for such an image similarity metric include simple $L_p$ norms (including $p=0$),
or more sophisticated metrics, like entropy based metrics \citep{umbrella}.
While no metric exists which reliable simulates human perception of image similarity, and $L_p$ norms may be insufficient in this regard \citep{lpnorms}, they still give us information about the robustness of the neural networks.
Here we will examine adversarial perturbations optimized towards a small $L_2$ norm, since it is often used in discussion about adversarial attacks and is a good compromise between the other commonly used, but more specific scenario targeting $L_0$ and $L_\infty$ norms.

Throughout this work $x \in [0,1]^n$ will denote the (flattened) input image and $\delta \in [0,1]^n$ the perturbation.
The label assigned to $x$ by the neural network will be referred to as $C(x)$, while the correct label is $C^*(x)$.\\
The networks output as logits is $Z(X)$ and the output interpretable as probabilities is $F(x)$.
This means, in the case of the ConvNet we have $F(x) = \mathrm{softmax}(Z(x))$ and in the case of the CapsNet $Z(x) = \mathrm{arctanh}(2F(x) - 1)$. \\
Furthermore, any gradient in the following section refers to the gradient with respect to the input image $x$ and not the network's parameters.

\subsection{Carlini-Wagner Attack}

\citet{carlini} extensively evaluated a wide range of adversarial attacks.
In particular they developed the following targeted white-box attack which will be referred to as the Carlini-Wagner (CW) attack.

Starting with the general problem for targeted attacks for a given label t
\begin{equation}
\begin{aligned}
& \minimize_{\delta} && \norm{\delta} \\
& \text{subject to} && C(x + \delta) = t \\
& && x + \delta \in [0,1]^n .
\end{aligned}
\end{equation}
Using Lagrangian relaxation, this is transformed to the following problem
\begin{equation}
\label{eq:cwproblem}
\begin{aligned}
& \minimize_{\delta} && \norm{\delta} + c \cdot f(x + \delta)\\
& \text{subject to} && x + \delta \in [0,1]^n ,
\end{aligned}
\end{equation}
where $c > 0$ is a suitable chosen constant and $f$ is an \emph{objective function}, i.e.\ a function with the property, that
$f(x + \delta) \leq 0$ if and only if $C(x + \delta) = t$.
There are many possible objective functions, but a particularly good choice is
\begin{equation}
\label{eq:cwobjective}
f(x') = \max \Set{ \max \Set{Z(x')_i | i \neq t } - Z(x')_t, -\kappa } .
\end{equation}%

The \emph{confidence parameter} $\kappa$ determines with how much confidence the adversarial example is classified as the target class, i.e.\ how close to the decision boundary it is.
The optimal value for $c$ in \Cref{eq:cwproblem} is the smallest value, that results in an adversarial example and can be found using a binary search.

This box-constrained problem can further be simplified by introducing the variable $w \in \mathbb{R}^n$ and setting $\delta_i = \frac{1}{2} (\tanh(w_i) + 1)$.
This substitution yields an unconstrained problem, which can be solved with various popular methods. We used the Adam optimizer \citep{adam}.

The Carlini-Wagner attack is generally a quite strong attack and leads to almost undetectable adversarial examples.
In particular, it can find adversarial examples even when defensive techniques like distillation are utilized \citep{carlini}.
However, due to the binary search it can often be rather slow.

\subsection{Boundary Attack}

The boundary attack is a black-box attack proposed by \citet{boundary}. It is furthermore a \emph{decision based} attack. This means, not only is no knowledge of the model architecture or the learned weights required, but also the output scores of the network are hidden from the attack. Only the final decision $C(x)$ can be used to construct an adversarial example.
Unlike other attacks, the boundary attack does not start with the original image and modifies it, but instead starts with a (possible random) misclassified image and changes it to resemble the original image.

The boundary attack iteratively generates adversarial examples $x^{(k)}$ (and therefore perturbations $\delta^{(k)} = x^{(k)} - x$) using a random walk. At the beginning it samples $x^{(0)} \sim \mathcal{U}(0,1)^n$ until $C(x^{(0)}) \neq C^*(x)$, i.e.\ $x^{(0)}$ is an adversarial example for $x$.
The next iteration should then fulfill the following criteria to ensure that the sequence $(x^{(k)})_{k \in \mathbb{N}}$ has a limit point on the decision boundary:

\begin{enumerate}
	\item The new image $x^{(k+1)}$ is in the range of a valid image:
	\begin{equation*}
	x^{(k+1)}\in [0,1]^n
	\end{equation*}
	
	\item The size of the random step is proportional to the size of the perturbation $\delta^{(k)}$ for a parameter $\gamma > 0$:
	\begin{equation*}
	\norm{\delta^{(k+1)} - \delta^{(k)}} = \gamma \cdot \norm{\delta^{(k)}}
	\end{equation*}
	
	\item The distance to the original is reduced with a factor $0<\nu<1$:
	\begin{equation*}
	\norm{\delta^{(k+1)}} = \nu \cdot \norm{\delta^{(k)}}
	\end{equation*}
\end{enumerate}

While these conditions are difficult to meet exactly, they are approximated by sampling from a normal distribution orthogonal to $\delta^{(k)}$ and making a step with size $\gamma$ in this direction.
If this point is still adversarial, a step towards $x$ is made with size $\nu$.
The parameters $\gamma$ and $\nu$ are adapted dynamically, similarly to Trust Region \citep{trustregion} methods.
A moving average of the success rate of the orthogonal of the direct step is kept and the parameters are increased or decrease, if this average deviates too much from beforehand chosen optimal value.

The success rate for the orthogonal step should be close to $0.5$, the idea being that if $x^{(k)}$ is close to the decision boundary, which is locally linear almost everywhere for most types of neural networks, then approximately $50\%$ of random orthogonal steps should pass the decision boundary. 
The target success rate for the direct step is set to a value around $0.25$, which empirically leads to good adversarial examples.
The algorithm terminates once $\nu$ is close to zero.

\subsection{DeepFool Attack}

\citet{deepfool} developed the untargeted white-box attack \emph{DeepFool}.
The authors calculate minimal adversarial perturbations for a linear classifier by projecting the original image to the nearest decision boundary.
By approximating the network with its first order Taylor polynomial, these calculations can be applied to the nonlinear classifier.
In detail, the perturbation is initialized as $\delta^{(0)} = 0$.
For each class label $i \neq C^*(x)$ the distance to the decision boundary is estimated with

\begin{equation}
%\begin{aligned}
l_i = 
\frac{
	\abs{Z_i(x + \delta^{(k)}) - Z_{C^*(x)}(x + \delta^{(k)})}
}
{
	\norm{\nabla Z_i(x + \delta^{(k)}) - \nabla Z_{C^*(x)}(x + \delta^{(k)})}
}
%\end{aligned}
\end{equation}

\todo{How is this similar to Carlini-Wagner?}

and $\delta^{(k+1)}$ is the projection to the nearest boundary under the distance approximation from above.
Our version of the DeepFool attack slightly differs from the original version in insofar that we clip $\delta^{(k)}$ to obtain $x + \delta^{(k)} \in [0,1]^n$, and that we restrict $\norm{\delta^{(k)}}$ at each step, which yields better adversarial examples in our experience.

\subsection{Universal Adversarial Perturbations}

The concept of universal adversarial perturbations was proposed by \citet{universal} and refers to a single perturbation vector $\delta \in \mathbb{R}^n$ such that $C(x + \delta) \neq C^*(x)$ for many different $x$ sampled from the input image distribution.
To do this, we use following variation of their algorithm:

As long as the accuracy on the test set is above a previously chosen threshold, repeat these steps:
\begin{enumerate}
	\item Initialize $\delta^{(0)} \gets 0$.
	\item Sample a batch $X^{(k)} = \{x_1^{(k)}, ..., x_N^{(k)}\}$ of images with $\forall x \in\ X^{(k)}:  C(x + \delta^{(k)}) = C^*(x)$.
	\item For each $x_i^{(k)}$ compute a perturbation $\delta_i^{(k+1)}$ using FGSM \citep{fgsm}.
	\item Update the perturbation: $\delta^{(k+1)} \gets \delta^{(k)} + \frac{1}{N} \sum\limits_{i=0}^N \delta_i^{(k+1)}$
\end{enumerate}

This differs insofar from the original algorithm described by \citet{universal} that in step 3 approximate perturbations for a whole batch is computed, while the computed an optimal perturbation for each $x_i$ using their previous work DeepFool.

In principle many other attacks instead of FGSM can be used, however we found that for our applications FGSM reached adequate results compared to much slower methods.
Since FGSM uses the gradient of the network, FGSM and therefore this algorithm for computing universal adversarial perturbations, are white-box attacks.

% !TeX root = ba.tex

\section{Experiments}

\subsection{Results}

\todo{Nicer looking tables}
\begin{table*}[h]
%	\begin{minipage}{.45\linewidth}
		\caption{Average perturbation norms for each attack and architecture.}
		\vskip 0.15in
		\centering\scalebox{1.0}{
			\begin{tabular}{llcccc}
				\toprule
				Attack & Network       & MNIST & Fashion & SVHN & CIFAR10  \\
				\midrule
				\multirow{2}{*}{CW} & ConvNet & {$1.40$} & $0.51$ & $0.67$ & $0.37$ \\
				& CapsNet            & $1.82$ & {$0.50$} & {$0.60$} & {$0.23$} \\
				\midrule
				\multirow{2}{*}{Boundary} & ConvNet & {$3.07$} & $1.24$ & $2.42$ & $1.38$ \\
				& CapsNet            & $3.26$ & {$0.93$} & {$1.88$} & {$0.72$} \\
				\midrule
				\multirow{2}{*}{DeepFool} & ConvNet & {$1.07$} & {$0.31$} & {$0.41$} & $0.23$ \\
				& CapsNet           & $2.02$ & $0.55$ & $0.80$ & {$0.16$} \\
				\midrule
				\multirow{2}{*}{Universal} & ConvNet & {$6.71$} & {$2.61$} & {$2.46$} & {$2.45$} \\
				& CapsNet           & $11.45$ & $5.31$ & $8.59$ & $2.70$ \\
				\bottomrule\\
		\end{tabular}}
		\label{tab:norms}
		
%	\end{minipage}\hspace{0.8cm}
%	\begin{minipage}{.45\linewidth}
		\caption{Fooling rates of adversarial examples calculated for a CapsNet and evaluated on a ConvNet and vice versa. For the universal attack we report the accuracy on the whole test set.}
		\vskip 0.15in
		\centering\scalebox{1.0}{
			\begin{tabular}{llcccc}
				\toprule
				Attack & Network       & MNIST & Fashion & SVHN & CIFAR10  \\
				\midrule
				\multirow{2}{*}{CW} & ConvNet & $0.8\%$ & $1.2\%$ & $2.8\%$ & $2.4\%$ \\
				& CapsNet            & $2.0\%$ & $2.0\%$ & $3.8\%$ & $2.0\%$ \\
				\midrule
				\multirow{2}{*}{Boundary} & ConvNet & $8.8\%$ & $9.5\%$ & $10.5\%$ & $13.4\%$ \\
				& CapsNet            & $14.2\%$ & $14.6\%$ & $12.9\%$ & $26.1\%$ \\
				\midrule
				\multirow{2}{*}{DeepFool} & ConvNet & $4.3\%$ & $8.5\%$ & $13.5\%$ & $11.8\%$ \\
				& CapsNet           & $0.9\%$ & $10.9\%$ & $10.8\%$ & $14.1\%$ \\
				\midrule
				\multirow{2}{*}{Universal} & ConvNet & $4.9\%$ & $20.4\%$ & $35.0\%$ & $25.9\%$ \\
				& CapsNet           & $38.2\%$ & $25.7\%$ & $53.4\%$ & $47.2\%$ \\
				\bottomrule\\
		\end{tabular}}
		\label{tab:attacks}%\end{minipage}
\end{table*}

%%%%%%%%%%%%%%%%%%%%%%%%%%%%%%%%%%%%%%%%%%%%%%%%%%%%%%%%%%%%%%%%%%%%%%%%
%%%% ENDE TEXTTEIL %%%%%%%%%%%%%%%%%%%%%%%%%%%%%%%%%%%%%%%%%%%%%%%%%%%%%
%%%%%%%%%%%%%%%%%%%%%%%%%%%%%%%%%%%%%%%%%%%%%%%%%%%%%%%%%%%%%%%%%%%%%%%%

\clearpage

% Entfernen Sie das Kommentar aus der nachfolgenden Zeile, falls Sie einen Anhang in der Arbeit verwenden wollen. Beachten Sie, dass Sie sich im Verlauf der Arbeit mit \ref{...} (z.B. \ref{anhang:zusatz1}) auf den Anhang beziehen.
%\newpage
\appendix
\section{Anhang}

\subsection*{Zusatzteil 1} \label{anhang:zusatz1}

Dies ist ein Anhang.

\clearpage

%\bibliography{references}
%\bibliographystyle{alphadin}
\DeclareNameAlias{sortname}{first-last}
\printbibliography[heading=bibintoc, title=References]
%\vspace*{\fill}

\clearpage

\listoffigures

\listoftables

%\pagebreak

%\printindex
\listoftodos
\end{document}
